% Part: sets-functions-relations
% Chapter: sets
% Section: pairs-and-products

\documentclass[../../../include/open-logic-section]{subfiles}

\begin{document}

\olfileid{sfr}{set}{pai}
\olsection{జతలు, క్రమిత బహుళకాలు, కార్టీజియన్ లబ్ధాలు}

\begin{explain}
సమితులలోని మూలకాలకు క్రమం ఉండదని సమితుల సమానత్వ సూత్రం చెబుతుంది.
కనుక క్రమాన్ని సూచించాలంటే $\tuple{x, y}$ అనే \emph{క్రమయుగ్మాలను}
వాడతాం. $\{x, y\}$ అనే అక్రమిత జతలో క్రమం ముఖ్యం కాదు:
$\{x, y\} = \{y, x\}$. క్రమయుగ్మంలో మాత్రం క్రమం ముఖ్యం:
$x \neq y$ అయితే $\tuple{x, y} \neq \tuple{y, x}$.

సమితి సిద్ధాంతంలో క్రమయుగ్మాలను ఎలా పరిగణించాలి? రెండు క్రమయుగ్మాలలో
మొదటి మూలకం ఒకటే అయి, రెండవ మూలకం కూడా ఒకటే అయినప్పుడే ఆ జతలు
సమానమవ్వాలన్న భావనను కాపాడాలి. అంటే:
\[
  \tuple{a, b}= \tuple{c, d}\text{ కావడానికి అవసరమైన, సరిపడే షరతులు }a = c \text{ మరియు }b=d.
\]
సమితి సిద్ధాంతంలో వీనర్--కురటోవ్‌స్కీ నిర్వచనాన్ని వాడి క్రమయుగ్మాలను
నిర్వచించవచ్చు.
\end{explain}

\begin{defn}[క్రమయుగ్మం]\ollabel{wienerkuratowski}
	$\tuple{a, b} = \{\{a\}, \{a, b\}\}$.
\end{defn}

\begin{prob}
	\olref[sfr][set][pai]{wienerkuratowski}ను ఉపయోగించి,
	$\tuple{a, b}= \tuple{c, d}$ కావడానికి అవసరమైన మరియు సరిపడే
	షరతులు $a = c$, $b=d$ రెండూ నెరవేరడమేనని నిరూపించండి.
\end{prob}

\begin{explain}
క్రమయుగ్మం నిర్వచనాన్ని నిర్ణయించాక, దానితో మరిన్ని సమితులను
నిర్వచించవచ్చు. ఉదాహరణకు కొన్నిసార్లు రెండుకంటే ఎక్కువ వస్తువుల
క్రమాలు అవసరమవుతాయి: $\tuple{x, y, z}$ వంటి \emph{త్రికాలు},
$\tuple{x, y, z, u}$ వంటి \emph{చతుష్కాలు}, మొదలైనవి.
మొదటి మూలకమే ఒక క్రమయుగ్మంగా ఉన్న ప్రత్యేక క్రమయుగ్మంగా త్రికాన్ని
పరిగణించవచ్చు: $\tuple{x, y, z}$ అనేది $\tuple{\tuple{x, y},z}$.
చతుష్కానికీ ఇదే విధంగా: $\tuple{x,y,z,u}$ అనేది
$\tuple{\tuple{\tuple{x,y},z},u}$. ఇలాగే కొనసాగించవచ్చు.
సాధారణంగా $\tuple{x_1, \dots, x_n}$లను \emph{క్రమిత $n$-బహుళకాలు}
అంటాం.

క్రమయుగ్మాలతో లేదా ఇతర క్రమిత $n$-బహుళకాలతో ఏర్పడే కొన్ని సమితులు
మనకు ఉపయోగపడతాయి.
\end{explain}

\begin{defn}[కార్టీజియన్ లబ్ధం]
$A$, $B$ అనే సమితులు ఇచ్చినప్పుడు, వాటి \emph{కార్టీజియన్ లబ్ధం}
$A \times B$ను ఇలా నిర్వచిస్తాం:
\[
  A \times B = \Setabs{\tuple{x, y}}{x \in A \text{ మరియు } y \in B}.
\]
\end{defn}

\begin{ex}
$A = \{0, 1\}$, $B = \{1, a, b\}$ అయితే, వాటి లబ్ధం
\[
A \times B = \{ \tuple{0, 1}, \tuple{0, a}, \tuple{0, b},
    \tuple{1, 1}, \tuple{1, a}, \tuple{1, b} \}.
\]
\end{ex}

\begin{ex}
$A$ ఒక సమితి అయితే, $A$తో దాని లబ్ధం $A \times A$ను $A^2$
అని కూడా రాస్తాం. ఇది $x, y \in A$ అయ్యే $\tuple{x, y}$ అనే
జతలన్నింటి సమితి; ఇక్కడ \emph{అన్ని} జతలనూ చేర్చాలి.
$\tuple{x, y, z}$ అనే త్రికాలన్నింటి సమితి $A^3$; ఇదే విధంగా
కొనసాగించవచ్చు. పునరావృత్త నిర్వచనాన్ని ఇలా ఇవ్వవచ్చు:
\begin{align*}
  A^1 & = A\\
  A^{k+1} & = A^k \times A
\end{align*}
\end{ex}

\begin{prob}
$\{1, 2, 3\}^3$లోని \tetoken{మూలకాలను}{!!{element}s} అన్నింటినీ జాబితాగా రాయండి.
\end{prob}

\begin{prop}\ollabel{cardnmprod}
$A$లో $n$ \tetoken{మూలకాలు}{!!{element}s}, $B$లో $m$ \tetoken{మూలకాలు}{!!{element}s} ఉంటే,
$A \times B$లో $n\cdot m$ మూలకాలు ఉంటాయి.
\end{prop}

\begin{proof}
$A$లోని ప్రతి \tetoken{మూలకం}{!!{element}} $x$కూ, $\tuple{x, y} \in A \times B$
రూపంలో $m$ \tetoken{మూలకాలు}{!!{element}s} ఉంటాయి. $B_x = \Setabs{\tuple{x, y}}{y
  \in B}$ అనుకుందాం. $x_1 \neq x_2$ అయినప్పుడల్లా
$\tuple{x_1, y} \neq \tuple{x_2, y}$ కనుక
$B_{x_1} \cap B_{x_2} = \emptyset$.
అయితే $A = \{x_1, \dots, x_n\}$ అయితే,
$A \times B = B_{x_1} \cup \dots \cup B_{x_n}$.
కాబట్టి అందులో $n\cdot m$ \tetoken{మూలకాలు}{!!{element}s} ఉంటాయి.

దీనిని చూడటానికి $A \times B$లోని \tetoken{మూలకాలను}{!!{element}s}
అడ్డువరుసలు, నిలువువరుసలుగా అమర్చుదాం:
\[
\begin{array}{rcccc}
  B_{x_1} = & \{\tuple{x_1, y_1} & \tuple{x_1, y_2} & \dots & \tuple{x_1, y_m}\}\\
  B_{x_2} = & \{\tuple{x_2, y_1} & \tuple{x_2, y_2} & \dots & \tuple{x_2, y_m}\}\\
  \vdots & & \vdots\\
  B_{x_n} = & \{\tuple{x_n, y_1} & \tuple{x_n, y_2} & \dots & \tuple{x_n, y_m}\}
\end{array}
\]
$x_i$లన్నీ వేర్వేరు, $y_j$లన్నీ వేర్వేరు కనుక ఈ అమరికలో
ఏ రెండు జతలూ ఒకేలా ఉండవు. మొత్తం $n\cdot m$ జతలు ఉంటాయి.
\end{proof}

\begin{prob}
$k$పై ఆగమనంతో, ప్రతి $k \ge 1$కూ ఈ విషయాన్ని చూపండి:
$A$లో $n$ \tetoken{మూలకాలు}{!!{element}s} ఉంటే, $A^k$లో $n^k$ \tetoken{మూలకాలు}{!!{element}s} ఉంటాయి.
\end{prob}

\begin{ex}
$A$ ఒక సమితి అయితే, $A$పై ఒక \emph{పదం} అంటే $A$లోని
\tetoken{మూలకాల}{!!{element}s} ఏదైనా క్రమం. అటువంటి క్రమాన్ని $A$లోని
\tetoken{మూలకాల}{!!{element}s} $n$-బహుళకంగా పరిగణించవచ్చు.
ఉదాహరణకు $A = \{a, b, c\}$ అయితే, ``$bac$'' అనే క్రమాన్ని
$\tuple{b, a, c}$ అనే త్రికంగా పరిగణించవచ్చు.
పదాలు, అంటే సంకేతాల క్రమాలు, కంప్యూటర్ శాస్త్రంలో అత్యంత ముఖ్యమైనవి.
ఒక ఆనవాయితీ ప్రకారం $A$లోని \tetoken{మూలకాలను}{!!{element}s}
$1$ పొడవు గల క్రమాలుగా, $\emptyset$ను $0$ పొడవు గల క్రమంగా
పరిగణిస్తాం. అప్పుడు $A$పై గల \emph{అన్ని} పదాల సమితి
\[
A^* = \{\emptyset\} \cup A \cup A^2 \cup A^3 \cup \dots
\]
\end{ex}

\end{document}
